%% abstract.tex: abstract in PT and EN  (FEUP regulations)
%% -------------------------------------------------------
% TODO write abstract
\chapter*{Resumo}
Sistemas de Produção Ciber-Físicos \acs{CPPS} place increasing demands on the software tools used to design and operate distributed industrial control applications. The \textit{IEC 61499} standard provides a structured and event-driven function block model well suited to these environments, yet existing development tools remain predominantly desktop-based and poorly adapted to collaborative or remote workflows. This dissertation investigates whether modern web technologies can meaningfully address these limitations and proposes \textit{PTERO}, a remote web-based \ac{IDE} for the modelling, deployment and monitoring of \textit{IEC 61499}-compliant applications.

The developed architecture shifts heavy orchestration workloads from local workstations to a centralised middleware server. To resolve the collaborative friction inherent to traditional file-based workflows, the platform incorporates a collaboration layer driven by \acp{CRDT}, facilitating concurrent multi-user visual modelling.

The system was validated through performance benchmarks against \textit{4diac IDE}. The evaluation demonstrated that \textit{PTERO} substantially reduced configuration overhead and improved the collaboration capabilities of the system. Furthermore, a usability study confirmed the participants' perception of these improvements. Ultimately, this research proves that migrating to a web-based engineering paradigm successfully modernizes distributed industrial automation.

\acresetall %% to reset the acronym usage

\chapter*{Abstract}
\acp{CPPS} place increasing demands on the software tools used to design and operate distributed industrial control applications. The \textit{IEC 61499} standard provides a structured and event-driven function block model well suited to these environments, yet existing development tools remain predominantly desktop-based and poorly adapted to collaborative or remote workflows. This dissertation investigates whether modern web technologies can meaningfully address these limitations and proposes \textit{PTERO}, a remote web-based \ac{IDE} for the modelling, deployment and monitoring of \textit{IEC 61499}-compliant applications.

The developed architecture shifts heavy orchestration workloads from local workstations to a centralised middleware server. To resolve the collaborative friction inherent to traditional file-based workflows, the platform incorporates a collaboration layer driven by \acp{CRDT}, facilitating concurrent multi-user visual modelling.

The system was validated through performance benchmarks against \textit{4diac IDE}. The evaluation demonstrated that \textit{PTERO} substantially reduced configuration overhead and improved the collaboration capabilities of the system. Furthermore, a usability study confirmed the participants' perception of these improvements. Ultimately, this research proves that migrating to a web-based engineering paradigm successfully modernizes distributed industrial automation.

\acresetall %% to reset the acronym usage
